% Blueprint content.
%
% Formalized nodes are written in Lean, in `blueprint/CoarseBlueprint.lean`, and inserted here with
% \inputleannode{label}; `lake build :blueprint` must run first. Nodes not formalized yet are
% written in LaTeX in the chapter files, without \notready. A Lean node whose statement awaits
% review is tagged `notReady := true`, which LeanArchitect renders as \notready.
%
% Each chapter is one milestone of the project. The list below names, for each, the key
% definitions and results; it says the same as the milestone's row in the roadmap and the
% chapter's introduction.

% Makes \inputleannode and \inputleanmodule available.
\input{../../.lake/build/blueprint/library/CoarseBlueprint}

This blueprint formalizes parts of coarse geometry and geometric group theory in Lean~4, in a
form that can move into Mathlib. Work is organized in \emph{milestones}, groups of definitions
and theorems chosen one at a time, with one chapter each. A chapter contains its main
statements, the definitions that Mathlib does not have yet and the statements need, then the
results their proofs need. Definitions and statements are formalized first, with the proofs left
open, and proved later; so a node may have its statement formalized long before its proof.

Each chapter follows a reference text. The default is C.~L\"oh, \emph{Geometric Group Theory:
An Introduction}, Springer, 2017; item numbers without a book refer to it. Hyperbolic spaces
also follow M.~R.~Bridson and A.~Haefliger, \emph{Metric Spaces of Non-Positive Curvature},
Springer, 1999. Coarse spaces of bounded geometry follow J.~Roe, \emph{Lectures on
Coarse Geometry}, American Mathematical Society, 2003.

The milestones, one per chapter, in order:
\begin{itemize}
  \item M1, \emph{Quasi-isometries} (Chapter~\ref{chap:quasi-isometry}): coarse Lipschitz and
        coarse antilipschitz maps, maps at finite distance, quasi-isometric embeddings,
        quasi-isometries and quasi-dense range; quasi-isometries are the quasi-isometric
        embeddings with quasi-dense range (Proposition 5.1.10), bilipschitz maps and coarse
        Lipschitz maps with a coarse Lipschitz quasi-inverse are quasi-isometries, and
        quasi-isometries compose.
  \item M2, \emph{Word metrics} (Chapter~\ref{chap:word-metric}): generating families, word
        length, word metrics and Cayley graphs, and a group with the word metric of a generating
        family; the Cayley-graph distance is the word metric, maps with bounded steps along the
        generators are Lipschitz for the word metric, word lengths for two finite generating
        families bound each other, and balls are finite.
  \item M3, \emph{Coarsely geodesic spaces and quasi-geodesics}
        (Chapter~\ref{chap:coarsely-geodesic}): the Rips graph, coarsely connected and coarsely
        geodesic spaces, quasi-geodesics; maps with bounded steps on coarsely geodesic spaces are
        coarse Lipschitz, normed spaces, geodesic spaces and word metrics are coarsely geodesic,
        and the coarsely geodesic spaces are L\"oh's quasi-geodesic spaces.
  \item M4, \emph{The \v{S}varc--Milnor lemma} (Chapter~\ref{chap:svarc-milnor}): metrically
        proper and cobounded actions; properly discontinuous actions, actions of finite groups and
        restrictions of metrically proper actions to subgroups are metrically proper, cocompact
        actions are cobounded, subgroups of finite index act coboundedly, and a group acts on
        itself with a word metric isometrically and metrically properly; a group acting
        isometrically, metrically properly and coboundedly on a coarsely geodesic space is
        finitely generated, and its orbit maps are quasi-isometries for the word metric of any
        finite generating family (Proposition 5.4.1); the topological form (Corollary 5.4.2);
        corollaries: independence of the word metric from the finite generating family up to
        quasi-isometry (Proposition 5.2.5), finite-index subgroups, quotients by finite normal
        subgroups (Corollary 5.4.7), and $\Z^n \sim_{QI} \R^n$ (Example 5.2.8).
  \item M5, \emph{Coarse spaces} (Chapter~\ref{chap:coarse-spaces}): the coarse-space library
        the project contains: coarse spaces, coarsely close maps, coarse maps, coarse
        equivalences and the metric coarse structure; the controlled sets of a pseudometric
        space, and a coarse \v{S}varc--Milnor lemma.
  \item M6, \emph{Uniformly discrete spaces and bounded geometry}
        (Chapter~\ref{chap:bounded-geometry}): uniformly discrete spaces and bounded geometry
        for metric spaces (Definition 5.6.11), bounded geometry and uniform local finiteness for
        coarse spaces (Roe, \S3.1); word metrics are uniformly discrete with bounded geometry,
        coarse bounds on uniformly discrete spaces are Lipschitz bounds (Proposition 5.2.9), and
        the metric and coarse notions agree on metric spaces.
  \item M7, \emph{Definitions of hyperbolicity} (Chapter~\ref{chap:hyperbolic}): Gromov's
        four-point condition, geodesic spaces and slim geodesic triangles,
        quasi-hyperbolic spaces (Definition 7.2.6), hyperbolic groups; the four-point condition
        and slim triangles agree on geodesic spaces (Exercise 7.E.10), quasi-hyperbolicity is
        invariant under quasi-isometric embeddings (Proposition 7.2.9), slim triangles and
        quasi-hyperbolicity agree on geodesic spaces (Theorem 7.2.10), Bridson--Haefliger's
        Theorem III.H.1.9, the four-point condition and quasi-hyperbolicity agree on word
        metrics, L\"oh's definition of hyperbolic groups (Definition 7.3.1), and the
        independence of hyperbolicity of a group from its generating family.
  \item M8, \emph{Geometric realisation and stability of quasi-geodesics}
        (Chapter~\ref{chap:stability}): the geometric realisation of a graph; its vertices sit
        isometrically and $\frac12$-densely and it is geodesic, the four-point condition passes
        from a quasi-dense isometric subspace to the whole space, and quasi-geodesics in spaces
        with slim triangles are stable (Bridson--Haefliger, Theorem III.H.1.7).
  \item M9, \emph{The Gromov boundary} (Chapter~\ref{chap:gromov-boundary}): the Gromov
        boundary of a hyperbolic space; not planned yet.
\end{itemize}

\chapter{Quasi-isometries}
\label{chap:quasi-isometry}

Milestone M1. All spaces are pseudometric spaces and all constants are in $\Rnn$. Elementary
facts about these notions (monotonicity in the constants, composition of the bounds, finite
distance being an equivalence relation, the calculations of L\"oh's Proposition 5.1.11) are not
nodes; the proof sketches below say where they are used.

\section{Coarse Lipschitz and coarse antilipschitz maps}

\inputleannode{Metric.CoarseLipschitzWith}
\inputleannode{Metric.CoarseAntilipschitzWith}

\section{Maps at finite distance}

\inputleannode{Metric.HasFiniteDist}

\section{Quasi-isometric embeddings}

\inputleannode{Metric.IsQuasiIsometricEmbeddingWith}
\inputleannode{Metric.IsQuasiIsometricEmbedding}
\inputleannode{Metric.IsQuasiIsometricEmbedding.of_coarseLipschitzWith_of_coarseAntilipschitzWith}

\section{Quasi-isometries}

\inputleannode{Metric.IsQuasiIsometry}
\inputleannode{Metric.QuasiIsometric}

\inputleannode{Metric.IsQuasiIsometry.comp}

\inputleannode{Metric.QuasiIsometric.refl}

\section{Quasi-dense range}

\inputleannode{Metric.QuasiDenseRange}
\inputleannode{Metric.isQuasiIsometry_iff_quasiDenseRange}

\chapter{Word metrics}
\label{chap:word-metric}

Milestone M3. $P$ is a generating family $\iota \to G$ of a group $G$, and $|g|_P$ the word
length of $g$ with respect to $P$. The formalized nodes of this chapter are copies of open
Mathlib pull requests (see \texttt{ATTRIBUTION.md}). Their elementary API (subadditivity, inverses,
word length zero, generators having length at most one) is used in the proofs below without
being a node.

\section{Generating families and word length}

\inputleannode{Group.Generators}
\inputleannode{Group.Generators.wordLength}
\inputleannode{Group.Generators.groupNorm}

\inputleannode{Group.Generators.cayleySimpleGraph}
\inputleannode{Group.Generators.dist_cayleySimpleGraph}

\section{Bounds}

\inputleannode{Group.Generators.dist_le_mul_wordLength}

\inputleannode{Group.Generators.wordLength_le_mul_wordLength}

\inputleannode{Group.Generators.finite_setOf_wordLength_le}

\chapter{Coarsely geodesic spaces}
\label{chap:coarsely-geodesic}

Milestone M2. An \emph{$r$-chain} from $x$ to $y$ is a sequence $x = x_0, \dots, x_n = y$ with
$d(x_i, x_{i+1}) \le r$; $n$ is its number of steps. A chain with a count is a map
$\mathrm{Fin}(n + 1) \to X$, and coarse connectedness uses the reflexive--transitive closure of
the relation $d(a, b) \le r$ (design decisions D6, D13).

\inputleannode{Metric.IsCoarselyConnectedWith}

\inputleannode{Metric.IsCoarselyGeodesicWith}

\inputleannode{NormedSpace.isCoarselyGeodesicWith}

\inputleannode{Group.Generators.isCoarselyGeodesicWith}

\chapter{The \v{S}varc--Milnor lemma}
\label{chap:svarc-milnor}

Milestone M5. $G$ acts isometrically on a nonempty pseudometric space $X$, $s \subseteq X$ is
bounded with $\bigcup_g g s = X$, $x_0 \in s$ and $r \ge 0$. Let $B$ be the closed
$r$-thickening of $s$, which is bounded. Both proofs use the set
\[ S = \{g \in G \mid g B \cap B \ne \emptyset\}, \]
the set in the definition of a metrically proper action, written in Lean with Mathlib's
\texttt{Metric.cthickening} and no new definition; it is L\"oh's generating set in Prop.~5.4.1.
The steps are stated with the word length and no metric on $G$ (design decision D8). The
hypotheses give a bounded $s$ whose translates cover $X$ but need not contain $x_0$; both proofs
replace $s$ by $s \cup \{x_0\}$, which is still bounded. Both proofs choose, for each $y \in X$, a
translate $\psi(y) s$ containing $y$; the map $\psi \colon X \to G$ is a quasi-inverse of the orbit
map, and $r$-close points have values of $\psi$ that differ by an element of $S$.

\inputleannode{Group.fg_of_isCoboundedSMul}

\inputleannode{Group.Generators.isQuasiIsometry_smul_const}

\section*{Topological form}

L\"oh states the lemma also with topological hypotheses (Corollary 5.4.2): a proper geodesic
space, and an action that is proper (for every compact $K$, only finitely many $g$ have $g K \cap
K \ne \emptyset$; Mathlib's \texttt{ProperlyDiscontinuousSMul}) and cocompact (the orbit space is
compact). These imply the hypotheses above
(Lemmas~\ref{isMetricallyProperSMul_of_properlyDiscontinuousSMul},
\ref{isCoboundedSMul_of_compactSpace_quotient} and
\ref{Metric.IsGeodesicSpace.isCoarselyGeodesicWith}).

\begin{corollary}[\v{S}varc--Milnor lemma, topological form, finite generation, L\"oh
  Cor.~5.4.2]
  \label{Group.fg_of_properlyDiscontinuousSMul}
  \uses{Metric.IsGeodesic}
  A group acting isometrically, properly discontinuously and cocompactly on a nonempty proper
  geodesic pseudometric space is finitely generated.
\end{corollary}
\begin{proof}
  \uses{Group.fg_of_isCoboundedSMul, isMetricallyProperSMul_of_properlyDiscontinuousSMul,
    isCoboundedSMul_of_compactSpace_quotient, Metric.IsGeodesicSpace.isCoarselyGeodesicWith}
  The action is metrically proper
  (Lemma~\ref{isMetricallyProperSMul_of_properlyDiscontinuousSMul}) and cobounded
  (Lemma~\ref{isCoboundedSMul_of_compactSpace_quotient}), and the space is coarsely geodesic
  (Lemma~\ref{Metric.IsGeodesicSpace.isCoarselyGeodesicWith}), hence coarsely connected, as it is
  nonempty. Apply Theorem~\ref{Group.fg_of_isCoboundedSMul}.
\end{proof}

\begin{corollary}[\v{S}varc--Milnor lemma, topological form, quasi-isometry, L\"oh Cor.~5.4.2]
  \label{Group.Generators.isQuasiIsometry_smul_const_of_properlyDiscontinuousSMul}
  \uses{Metric.IsGeodesic, Metric.IsQuasiIsometry, Group.Generators.WithWordMetric}
  Under the same hypotheses, every orbit map $g \mapsto g x_0$ is a quasi-isometry from $G$ with
  the word metric of any finite generating family.
\end{corollary}
\begin{proof}
  \uses{Group.Generators.isQuasiIsometry_smul_const,
    isMetricallyProperSMul_of_properlyDiscontinuousSMul, isCoboundedSMul_of_compactSpace_quotient,
    Metric.IsGeodesicSpace.isCoarselyGeodesicWith}
  As for Corollary~\ref{Group.fg_of_properlyDiscontinuousSMul}, with
  Theorem~\ref{Group.Generators.isQuasiIsometry_smul_const} in place of
  Theorem~\ref{Group.fg_of_isCoboundedSMul}.
\end{proof}

\chapter{Coarse spaces}
\label{chap:coarse-spaces}

The project contains a copy of S.~Ghobash's coarse-space library (see \texttt{ATTRIBUTION.md}).
It proves a \v{S}varc--Milnor statement as a coarse equivalence, without constants, word metric
or finite generation. Nothing yet connects it to the metric notions above.

\inputleannode{def:coarse-space}
\inputleannode{def:coarsely-close}
\inputleannode{def:coarse-map}
\inputleannode{def:coarse-equiv}
\inputleannode{def:metric-coarse-space}
\inputleannode{lem:coarsely-proper-smul}
\inputleannode{thm:coarse-svarc-milnor}

\section*{Candidate bridges}

These have no consumer yet; they become ready when a node uses them.

\begin{lemma}
  \label{lem:finite-dist-iff-coarsely-close}
  \notready
  \uses{def:finite-dist, def:coarsely-close, def:metric-coarse-space}
  Two maps into a pseudometric space have finite distance if and only if they are coarsely
  close for the metric coarse structure.
  \plannedlean{Metric.hasFiniteDist_iff_isCoarselyClose}
\end{lemma}
\begin{proof}
  Unfold the controlled sets of the metric coarse structure.
\end{proof}

\begin{lemma}
  \label{lem:qi-coarse-equiv}
  \notready
  \uses{def:qi, def:coarse-equiv, def:metric-coarse-space, lem:finite-dist-iff-coarsely-close}
  A quasi-isometry together with a quasi-inverse is a coarse equivalence for the metric coarse
  structures.
  \plannedlean{Metric.IsQuasiIsometry.toCoarseEquiv}
\end{lemma}
\begin{proof}
  Coarse Lipschitz maps are controlled, coarse antilipschitz maps are coarsely proper, and maps
  at finite distance are coarsely close.
\end{proof}

\chapter{Uniformly discrete spaces and bounded geometry}
\label{chap:bounded-geometry}

Milestone M8. A space has \emph{bounded geometry} if, at each scale, its balls are uniformly
small. There are two settings. For metric spaces we follow L\"oh (Definition 5.6.11), who
counts the points of closed balls and pairs bounded geometry with uniform discreteness. For
coarse spaces we follow J.~Roe, \emph{Lectures on Coarse Geometry}, AMS, 2003 (\S3.1), who
measures a set at the scale of a controlled set $E$, the \emph{gauge}, by its largest
$E$-separated subsets. A set $T$ is \emph{$E$-separated} if no two distinct points of $T$ are
$E$-related (\texttt{SetRel.IsSeparated} in Mathlib), and the \emph{$E$-ball} at $x$ is
$E_x = \{y \mid (y, x) \in E\}$ (\texttt{SetRel.ball}).

Uniform discreteness is a small-scale condition that coarse structures cannot see; Roe defines
it for metric spaces only (Example 3.14). On the coarse side its role is played by the
condition that the diagonal is a gauge. Uniformly discrete metric spaces are Mathlib's
\texttt{DiscreteUniformity} (design decision D22). The coarse counterpart and its name are an
open decision (D24); the coarse nodes follow its proposal.

\section*{Metric spaces}

\begin{definition}[Uniformly discrete space, L\"oh Def.~5.6.11]
  \label{DiscreteUniformity}
  A pseudometric space $X$ is \emph{uniformly discrete} if there is $\varepsilon > 0$ with
  $d(x, y) \ge \varepsilon$ for all $x \ne y$. In Lean this is Mathlib's class
  \texttt{DiscreteUniformity X}, which says that the diagonal is an entourage of the
  uniformity; for a pseudometric space the two agree, since the sets
  $\{(x, y) \mid d(x, y) < \varepsilon\}$, $\varepsilon > 0$, form a basis of the uniformity.
  \plannedlean{Metric.discreteUniformity_iff}
\end{definition}

\begin{definition}[Bounded geometry, L\"oh Def.~5.6.11]
  \label{Metric.HasBoundedGeometry}
  A pseudometric space $X$ has \emph{bounded geometry} if for every $r$ there is $K \in \N$
  such that every closed ball of radius $r$ in $X$ has at most $K$ points.
\end{definition}

L\"oh calls a metric space that is uniformly discrete and has bounded geometry a \emph{UDBG
space}; the conjunction gets no name here. Bounded geometry alone does not make a space
uniformly discrete: $\{n\} \cup \{n + 1/n\}$, $n \ge 1$, in $\R$ is a counterexample.

\begin{lemma}[Word metrics, cf.\ L\"oh after Def.~5.6.11]
  \label{Group.Generators.WithWordMetric.hasBoundedGeometry}
  \uses{Group.Generators.WithWordMetric, DiscreteUniformity, Metric.HasBoundedGeometry}
  A group with the word metric of a generating family is uniformly discrete. If the family is
  finite, the word metric has bounded geometry.
  \plannedlean{Group.Generators.WithWordMetric.instDiscreteUniformity}
\end{lemma}
\begin{proof}
  \uses{Group.Generators.finite_setOf_wordLength_le}
  Distances are word lengths, which are natural numbers, and only the identity has length
  $0$; so distinct elements are at distance at least $1$. Left multiplication by $g$ is an
  isometry taking the closed ball of radius $r$ about $1$ onto the one about $g$, so all closed
  balls of radius $r$ have as many points as $\{g \mid |g|_S \le \lfloor r \rfloor\}$, which is
  finite for a finite family (Lemma~\ref{Group.Generators.finite_setOf_wordLength_le}).
\end{proof}

\begin{lemma}[Coarse bounds on uniformly discrete spaces, cf.\ L\"oh Prop.~5.2.9]
  \label{Metric.CoarseLipschitzWith.exists_lipschitzWith}
  \uses{DiscreteUniformity, Metric.CoarseLipschitzWith, Metric.CoarseAntilipschitzWith}
  Let $f \colon X \to Y$ be a map of pseudometric spaces.
  \begin{enumerate}
    \item If $X$ is uniformly discrete and $f$ is coarse Lipschitz, then $f$ is Lipschitz.
    \item If $Y$ is uniformly discrete and $f$ is injective and coarse antilipschitz, then $f$ is
          antilipschitz.
  \end{enumerate}
  In particular an injective quasi-isometric embedding between uniformly discrete spaces is a
  bilipschitz embedding.
  \plannedlean{Metric.CoarseAntilipschitzWith.exists_antilipschitzWith}
\end{lemma}
\begin{proof}
  Let $\varepsilon > 0$ bound the distances between distinct points from below
  (Definition~\ref{DiscreteUniformity}). (1) If $d(f x, f y) \le L\,d(x, y) + C$, then for $x \ne
  y$, $C \le (C/\varepsilon)\,d(x, y)$, so $d(f x, f y) \le (L + C/\varepsilon)\,d(x, y)$; for $x =
  y$ both sides are $0$. (2) If $d(x, y) \le K\,d(f x, f y) + C$, then for $x \ne y$, $f x \ne f y$,
  so $C \le (C/\varepsilon)\,d(f x, f y)$, and $d(x, y) \le (K + C/\varepsilon)\,d(f x, f y)$.
\end{proof}

With Lemma~\ref{Group.Generators.WithWordMetric.hasBoundedGeometry}, this gives L\"oh's
Proposition 5.2.9: a bijective quasi-isometry between finitely generated groups with the word
metrics of finite generating families is a bilipschitz equivalence. Injectivity suffices, and the
Lipschitz bound needs neither injectivity nor a uniformly discrete codomain.

\section*{Coarse spaces}

\paragraph{Roe's definition.} Let $X$ be a coarse space. For $E \subseteq X \times X$ and
$S \subseteq X$, the \emph{$E$-capacity} $\operatorname{cap}_E(S) \in \N \cup \{\infty\}$ is the
largest size of an $E$-separated subset of $S$ (Roe, Definition 3.1). A \emph{gauge} is a
symmetric $E$ containing the diagonal $\Delta$. A set $F \subseteq X \times X$ is \emph{uniform}
with respect to a gauge $E$ if $\operatorname{cap}_E(F[E_x])$ and
$\operatorname{cap}_E(F^{-1}[E_x])$ are bounded uniformly in $x$, where
$F[A] = \{y \mid \exists a \in A,\ (y, a) \in F\}$ (Definition 3.5). The uniform sets form a
coarse structure $\mathcal{U}_E$, the \emph{universal bounded geometry structure} based on $E$
(Proposition 3.7, Definition 3.8), and $X$ has \emph{bounded geometry} if its coarse structure
is contained in $\mathcal{U}_E$ for some controlled gauge $E$ (Definition 3.9).

The definitions below state the same notion with four simplifications, listed after them with
the reason for each and the argument that nothing is lost. None of them is forced by Lean:
Roe's definition could be formalized as it stands. The first avoids proving a theorem before
the definition can be stated; the other three remove a construction or a hypothesis that the
definition does not need.

\begin{definition}[Bounded geometry with a gauge, cf.\ Roe Def.~3.5 and~3.9]
  \label{CoarseSpace.HasBoundedGeometryWith}
  \uses{CoarseSpace}
  Let $X$ be a coarse space and $E \subseteq X \times X$. The space has \emph{bounded geometry
  with gauge $E$} if for every controlled set $F$ there is $K \in \N$ such that for every
  $x \in X$, every $E$-separated subset of the ball $F_x$ has at most $K$ points.
\end{definition}

\begin{definition}[Bounded geometry, Roe Def.~3.9]
  \label{CoarseSpace.HasBoundedGeometry}
  \uses{CoarseSpace.HasBoundedGeometryWith}
  A coarse space has \emph{bounded geometry} if it has bounded geometry with some controlled
  gauge $E$.
\end{definition}

\subsection*{Differences from Roe's definition}

\begin{enumerate}
  \item \emph{No universal bounded geometry structure.} Roe says that the coarse structure is
        contained in $\mathcal{U}_E$. Stated literally in Lean, this needs $\mathcal{U}_E$ as
        a coarse structure (a \texttt{CoarseSpace}), and a coarse structure must be closed
        under products: that is Roe's Proposition 3.7, whose proof needs entropy as well as
        capacity (Lemma 3.6, Proposition 3.2(d)). The definition would then depend on a
        theorem. Containment only says that every controlled set is uniform with respect to
        $E$, and that is what Definition~\ref{CoarseSpace.HasBoundedGeometryWith} states, as a
        property of the coarse structure of $X$; no second coarse structure is built.
  \item \emph{No capacity function.} Roe measures sets by $\operatorname{cap}_E$, a supremum
        in $\N \cup \{\infty\}$. Here the definition says directly that every $E$-separated
        subset has at most $K$ points. Reason: the definition only ever compares a capacity
        with a bound, and a capacity function would need its own API (values in $\N \cup
        \{\infty\}$, suprema over subsets) with no other use yet; Mathlib has the metric
        version, \texttt{Metric.packingNumber}, but none for relations. Nothing is lost:
        $\operatorname{cap}_E(S) \le K$ exactly when every $E$-separated subset of $S$ has at
        most $K$ points (an infinite separated set has finite separated subsets of every
        size). If a later node needs capacities (growth, Roe's Proposition 3.15), the
        function can be added and this definition restated through it.
  \item \emph{No conditions on the gauge.} Roe's gauge is symmetric and contains the
        diagonal; here $E$ is any controlled set. Reason: the conditions are hypotheses every
        user would have to carry and check, and they do not change which sets are separated.
        A set is $E$-separated when no two \emph{distinct} points of it are $E$-related in
        either order (\texttt{SetRel.IsSeparated}); so adding $E^{-1}$ or $\Delta$ to $E$ does
        not change the $E$-separated sets. Hence if $E$ satisfies the definition, so does the
        controlled gauge $E \cup E^{-1} \cup \Delta$, with the same constants; and a gauge in
        Roe's sense is in particular a controlled set. The one place where Roe uses that a
        gauge contains $\Delta$ is item 4.
  \item \emph{Balls $F_x$ in place of $F[E_x]$ and $F^{-1}[E_x]$.} Roe bounds separated
        subsets of two sets for each controlled $F$; here, of one ball $F_x$. Reason: one
        condition instead of two, and no image of a set under a relation. Nothing is lost
        because the definition quantifies over \emph{all} controlled sets, which are closed
        under inverses and products: $F[E_x]$ is the ball at $x$ of $F \circ E$, and
        $F^{-1}[E_x]$ that of $F^{-1} \circ E$, both controlled when $E$ and $F$ are. So Roe's
        two conditions for $F$ are this definition's condition for $F \circ E$ and
        $F^{-1} \circ E$. Conversely $F_x \subseteq F[E_x]$, as $x \in E_x$ for a gauge.
        This step is valid only for the condition on all controlled sets at once, not for the
        uniformity of a single set $F$ (Roe's Definition 3.5 itself), which is why that
        notion is not defined here.
\end{enumerate}

Two further differences are not choices of this chapter. Roe's coarse structure is a
collection of controlled sets; Lean's \texttt{CoarseSpace} (Definition~\ref{CoarseSpace})
stores the filter of their complements, an equivalent encoding. And ``bounded uniformly in
$x$'' is written as a bound $K \in \N$ on cardinalities taken in $\N \cup \{\infty\}$, so that
infinite separated sets violate it. The equivalence with Roe's definition argued in items 1--4
is not formalized: it would need Roe's capacity, gauges and $\mathcal{U}_E$, which are
not defined here.

\begin{definition}[Uniformly locally finite coarse space, cf.\ Roe Ex.~2.44]
  \label{CoarseSpace.IsUniformlyLocallyFinite}
  \uses{CoarseSpace.HasBoundedGeometryWith}
  A coarse space is \emph{uniformly locally finite} if it has bounded geometry with the
  diagonal as gauge: for every controlled set $F$ there is $K \in \N$ such that every ball
  $F_x$ has at most $K$ points.
\end{definition}

Every set is separated for the diagonal, which gives the second form of the definition. The
diagonal is controlled, so a uniformly locally finite coarse space has bounded geometry.

\subsection*{Differences from Roe's Example 2.44}

Roe has no name for this condition. In Example 2.44 he defines, on a set $X$, the coarse
structure whose controlled sets $E$ have sections $E_x$ and $E^x = E^{-1}[\{x\}]$ of uniformly
bounded size, and Definition 3.8 identifies it with $\mathcal{U}_\Delta$. Here, as in item 1,
that coarse structure is not built: the definition says that every controlled set of the given
coarse structure satisfies Roe's condition, which is containment in $\mathcal{U}_\Delta$.
Only the sections $F_x$ are bounded, not $F^x$: the sections $F^x$ are the balls of the
controlled set $F^{-1}$ (item 4). The name ``uniformly locally finite'' is a proposal (design
decision D24), after the name of the same condition for graphs and simplicial complexes.

\section*{Comparison on metric spaces}

A pseudometric space carries the coarse structure in which a set of pairs is controlled when
the distance is bounded on it
(Definition~\ref{PseudoEMetricSpace.toCoarseSpace}).

\begin{lemma}[cf.\ Roe Ex.~3.14]
  \label{Metric.hasBoundedGeometry_iff_isUniformlyLocallyFinite}
  \uses{Metric.HasBoundedGeometry, CoarseSpace.IsUniformlyLocallyFinite,
    PseudoEMetricSpace.toCoarseSpace}
  A pseudometric space has bounded geometry if and only if its metric coarse structure is
  uniformly locally finite.
\end{lemma}
\begin{proof}
  \uses{Metric.isControlled_iff}
  The set $\{(x, y) \mid d(x, y) \le r\}$ is controlled, and its ball at $x$ is the closed ball
  of radius $r$ about $x$. Every controlled set is contained in one of these sets, so its balls
  are contained in closed balls of a fixed radius.
\end{proof}

\begin{theorem}[cf.\ Roe Ex.~3.13]
  \label{CoarseSpace.hasBoundedGeometry_iff_packingNumber_le}
  \uses{CoarseSpace.HasBoundedGeometry, PseudoEMetricSpace.toCoarseSpace}
  The metric coarse structure of a pseudometric space $X$ has bounded geometry if and only if
  there is $\varepsilon \ge 0$ such that for every $r$ there is $K \in \N$ with the following
  property: for every $x \in X$, every subset of the closed ball of radius $r$ about $x$ whose
  points are pairwise at distance greater than $\varepsilon$ has at most $K$ points. (The
  largest size of such a subset is Mathlib's packing number \texttt{Metric.packingNumber}.)
\end{theorem}
\begin{proof}
  \uses{Metric.isControlled_iff}
  For $\varepsilon \ge 0$ write $E_\varepsilon = \{(x, y) \mid d(x, y) \le \varepsilon\}$; a set
  is $E_\varepsilon$-separated exactly when its points are pairwise at distance greater than
  $\varepsilon$ (\texttt{Metric.isSeparated\_iff\_setRelIsSeparated}).
  ($\Rightarrow$) A controlled gauge $E$ lies in some $E_\varepsilon$, so every
  $E_\varepsilon$-separated set is $E$-separated; apply the definition to $F = E_r$, whose
  balls are the closed balls of radius $r$.
  ($\Leftarrow$) Take $E_\varepsilon$ as the gauge. A controlled $F$ lies in some $E_r$, so its
  balls lie in closed balls of radius $r$, and the bound for $r$ applies.
\end{proof}

Combining the two results with the remark after
Definition~\ref{CoarseSpace.IsUniformlyLocallyFinite}: a pseudometric space of bounded geometry
in L\"oh's sense has bounded geometry in Roe's sense (take $\varepsilon = 0$). The converse
fails, since bounded geometry in Roe's sense is invariant under coarse equivalence and L\"oh's
is not: $\R$ is coarsely equivalent to $\Z$.

\chapter{Hyperbolic spaces and groups}
\label{chap:hyperbolic}

Milestone M7. The definitions of hyperbolic spaces and groups below are copies of open Mathlib
pull requests. They use Gromov's four-point condition, which makes sense in every pseudometric
space. L\"oh (Definition 7.2.2) and Bridson--Haefliger (Definition III.H.1.1) define hyperbolicity
of a geodesic space by slim geodesic triangles instead. This chapter adds that definition and shows
that the two agree on geodesic spaces (L\"oh, Exercise 7.E.10; Bridson--Haefliger,
Proposition III.H.1.22).

Write $(x \cdot y)_w = \frac12\bigl(d(w, x) + d(w, y) - d(x, y)\bigr)$ for the Gromov product
(\texttt{Metric.gromovProduct w x y} in Mathlib; Bridson--Haefliger, Definition III.H.1.19). The
four-point condition with constant $\delta$ says
$(x \cdot z)_w \ge \min\{(x \cdot y)_w, (y \cdot z)_w\} - \delta$ for all $w, x, y, z$; this is
$(\delta)$-hyperbolicity in the sense of Bridson--Haefliger, Definition III.H.1.20.

Mathlib has no geodesics. Here a geodesic is an isometric map from an interval, as in L\"oh and
Bridson--Haefliger (design decision D17).
Chapter~\ref{chap:hyperbolic-generators} uses the comparison to show that hyperbolicity of a group
does not depend on the generating family.

\inputleannode{Metric.IsHyperbolicWith}
\inputleannode{Group.Generators.IsHyperbolicWith}

\inputleannode{Metric.IsGeodesic}

\inputleannode{Metric.HasSlimTrianglesWith}

L\"oh and Bridson--Haefliger call a geodesic space with $\delta$-slim triangles
$\delta$-hyperbolic. Here that name belongs to the four-point condition, and the constants of
the two notions differ.

Bridson--Haefliger prove Theorem~\ref{Metric.isHyperbolic_iff_exists_hasSlimTrianglesWith} through
thin triangles and the
insize of a triangle (Definition III.H.1.16, Propositions III.H.1.17 and III.H.1.22). The two
lemmas below go directly between slim triangles and the four-point condition instead
(design decision D19), so neither thin triangles nor insize is needed, and the constant is
$3\delta$ in both directions.

\inputleannode{Metric.HasSlimTrianglesWith.isHyperbolicWith}

\inputleannode{Metric.IsHyperbolicWith.hasSlimTrianglesWith}

\inputleannode{Metric.isHyperbolic_iff_exists_hasSlimTrianglesWith}

\chapter{Geometric realisation and stability of quasi-geodesics}
\label{chap:stability}

Milestone M8. This chapter provides the tools for the equivalences of
Chapter~\ref{chap:hyperbolic} that pass between exact distances and quasi-isometries: the
geometric realisation of a graph, a geodesic space in which the vertices sit isometrically and
$\frac12$-densely; the four-point condition on a space with a quasi-dense isometric copy of a
space satisfying it; and the stability of quasi-geodesics in spaces with slim triangles
(Bridson--Haefliger, Theorem III.H.1.7). With them,
Theorem~\ref{Group.Generators.isQuasiHyperbolic_iff} replaces the word metric of a group by the
realisation of its Cayley graph (design decisions D18 and D28).

How the realisation is encoded in Lean is an open decision (D26); its nodes are not implemented
until it is decided. Stability is proved with chains of points instead of rectifiable paths (design
decision D27).

\section*{Geometric realisation of a graph}

\begin{definition}[Geometric realisation, cf.\ Bridson--Haefliger Def.~I.1.9]
  \label{SimpleGraph.Realization}
  Let $\Gamma$ be a connected simple graph with vertex set $V$ and graph distance $d_\Gamma$. Its
  \emph{geometric realisation} $|\Gamma|$ is obtained from a copy of $[0, 1]$ for each edge by
  gluing the endpoints of the copy for the edge $\{u, u'\}$ to the vertices $u$ and $u'$. A point
  $p$ of $|\Gamma|$ is a vertex or lies inside an edge; a point at parameter $s$ on an edge from $u$
  to $u'$ has distance $s$ from $u$ and $1 - s$ from $u'$. The distance between $p$ on an edge with
  endpoints $u, u'$ and $q$ on an edge with endpoints $v, v'$ is the least of the numbers
  \[ d(p, a) + d_\Gamma(a, b) + d(b, q), \qquad a \in \{u, u'\},\ b \in \{v, v'\}, \]
  and, if $p$ and $q$ lie on the same edge, of their difference of parameters. Each vertex $v$ is a
  point of $|\Gamma|$, which defines the \emph{vertex map} $V \to |\Gamma|$.
  \plannedlean{SimpleGraph.Realization.vertex}
\end{definition}

\paragraph{Differences from Bridson--Haefliger.} Bridson--Haefliger glue copies of $[0, 1]$ for
the edges of any graph with edge lengths $\lambda(e) > 0$, and define the distance as the infimum of
the lengths of piecewise linear paths; that this is a geodesic metric when there are finitely many
edge lengths is left to the reader (I.1.9, fact 2). Here every edge has length $1$, and the
distance is given by the formula above, which for unit edges is the length of a shortest path: a
shortest path from $p$ to $q$ either stays in one edge or leaves the edge of $p$ through an
endpoint $a$, follows a shortest walk from $a$ to the endpoint $b$ of the edge of $q$ through
which it arrives, and ends inside that edge. Only unit edges are needed (Cayley graphs), and the formula
reduces every fact to a finite case analysis on top of the graph distance. Whether the
construction follows the formula or the quotient is design decision D26.

\begin{lemma}[Vertices of the realisation]
  \label{SimpleGraph.Realization.dist_vertex}
  \uses{SimpleGraph.Realization, Metric.QuasiDenseRange}
  The vertex map of a connected graph $\Gamma$ preserves distances: the distance in $|\Gamma|$
  between two vertices is their graph distance. Every point of $|\Gamma|$ is within $\frac12$ of a
  vertex; in particular the vertex map has quasi-dense range.
  \plannedlean{SimpleGraph.Realization.quasiDenseRange_vertex}
\end{lemma}
\begin{proof}
  A vertex is an endpoint of each edge it lies on, at distance $0$ from itself, so the formula for
  two vertices $u, v$ has the term $d_\Gamma(u, v)$. Every other term is at least
  $d_\Gamma(u, v)$: a term through endpoints $a, b$ is at least $d_\Gamma(u, a) + d_\Gamma(a, b) +
  d_\Gamma(b, v)$, since a point at distance $1$ from a vertex along an edge is the other endpoint,
  adjacent to it. A point at parameter $s$ on an edge from $u$ to $u'$ is within $\min\{s, 1 - s\}
  \le \frac12$ of $u$ or $u'$.
\end{proof}

\begin{lemma}[The realisation is geodesic, cf.\ Bridson--Haefliger I.1.9]
  \label{SimpleGraph.Realization.isGeodesicSpace}
  \uses{SimpleGraph.Realization, Metric.IsGeodesic}
  The geometric realisation of a connected graph is a geodesic space.
\end{lemma}
\begin{proof}
  If $p, q$ lie on one edge and the minimum is their difference of parameters, the segment of the
  edge between them is a geodesic. Otherwise take endpoints $a, b$ attaining the minimum and a walk
  $a = w_0, \dots, w_n = b$ with $n = d_\Gamma(a, b)$. Run from $p$ to $a$ along its edge, then
  along the edges $w_i w_{i+1}$ at unit speed, then from $b$ to $q$. This map $\gamma$ is defined on
  $[0, d(p, q)]$, $1$-Lipschitz on each piece and hence everywhere, and joins $p$ to $q$. A
  $1$-Lipschitz map $\gamma \colon [0, \ell] \to X$ with $d(\gamma(0), \gamma(\ell)) = \ell$ is
  isometric: for $s \le t$, $d(\gamma(s), \gamma(t)) \ge \ell - d(\gamma(0), \gamma(s)) -
  d(\gamma(t), \gamma(\ell)) \ge \ell - s - (\ell - t) = t - s$.
\end{proof}

\section*{The four-point condition and quasi-dense subspaces}

The four-point condition passes to subspaces with the same constant
(\texttt{Metric.IsHyperbolicWith.of\_isometry} in the copied pull request). It also passes from a
quasi-dense subspace to the whole space, with a larger constant; this is how the word metric and the
realisation of the Cayley graph are compared.

\begin{lemma}[Four-point condition from a quasi-dense subspace]
  \label{Metric.IsHyperbolicWith.of_isometry_of_quasiDenseRange}
  \uses{Metric.IsHyperbolicWith, Metric.QuasiDenseRange}
  Let $f \colon X \to Y$ be an isometric map of pseudometric spaces such that every point of $Y$ is
  within $C$ of a point of $f(X)$. If $X$ satisfies the four-point condition with constant
  $\delta$, then $Y$ satisfies it with constant $\delta + 6C$.
\end{lemma}
\begin{proof}
  Given $w, x, y, z \in Y$, choose $w', x', y', z' \in X$ whose images are within $C$ of them.
  Each distance between two of the four points changes by at most $2C$ when they are replaced by the
  images, so each Gromov product, half a sum of three distances, changes by at most $3C$. The
  four-point condition for $w', x', y', z'$ in $X$, transported by the isometry, gives $(x \cdot
  z)_w \ge (fx' \cdot fz')_{fw'} - 3C \ge \min\{(fx' \cdot fy')_{fw'}, (fy' \cdot fz')_{fw'}\} -
  \delta - 3C \ge \min\{(x \cdot y)_w, (y \cdot z)_w\} - \delta - 6C$.
\end{proof}

\section*{Stability of quasi-geodesics}

Quasi-geodesics are defined in Chapter~\ref{chap:coarsely-geodesic}
(Definition~\ref{Metric.IsQuasiGeodesicWith}).

\inputleannode{Metric.HasSlimTrianglesWith.exists_hausdorffDist_le}

\chapter{The Gromov boundary}
\label{chap:gromov-boundary}

Milestone M9. The Gromov boundary of a hyperbolic space (Chapter~\ref{chap:hyperbolic}). Its
definitions and results are not planned yet.

